\documentclass[openany,fontset=fandol]{ctexbook}
\usepackage{amsmath,amssymb}
\usepackage[unicode,hidelinks]{hyperref}
\setlength{\emergencystretch}{2em}
\raggedbottom
\DeclareMathOperator{\Spec}{Spec}
\DeclareMathOperator{\Ass}{Ass}
\DeclareMathOperator{\codim}{codim}
\DeclareMathOperator{\depth}{depth}
\DeclareMathOperator{\Supp}{Supp}
\DeclareMathOperator{\Gr}{Gr}
\DeclareMathOperator{\coprof}{coprof}
\begin{document}
\title{代数几何基础第五卷预备笔记\\例外超平面的余维}
\author{第100批\quad 法、英、中文三方校订}
\date{2026年11月24日}
\maketitle
\section*{范围与接续}
本批译校第8节标题、8.1、8.2及其末尾附记：法扫描物理30--31页和32页开头，对照英编辑稿物理21--22页。法30的划销旧稿单列，不与法31的正稿拼接。下一批从法32的8.3开始；1158页合订本未改。
\section*{第8节\quad 例外超平面集合的维数}
\textbf{8.1.}
在前面的各节，特别是第2、3节，我们曾给出一些断言：在对偶射影空间 $P^\vee$ 中，使超平面截面 $Y_\xi$ 具有某一性质 $\mathcal P$ 的参数 $\xi$ 构成可构造集，并包含泛点 $\eta$。换言之，对性质 $\mathcal P$ 而言的例外参数集 $Z_{\mathcal P}$ 是可构造且稀薄的，即其闭包的余维至少为 $1$。这里假设 $S=\Spec k$。

在某些情形下，可以为这个余维给出更强的下界，从而细化上述断言。\footnote{法稿使用表示上界的 \emph{majoration}，英稿也写作上界；而下一句实际用的是余维 $\ge 2$，故正文改作下界。若改说例外闭包的维数，才是更强的上界。}
这对若干问题是重要的。例如，若该余维至少为 $2$，则 $P^\vee$ 中一条充分一般的直线 $D$ 不与 $Z_{\mathcal P}$ 相遇。因而，若 $k$ 是无限域，就存在由 $D$ 的几何点 $\xi$ 参数化的超平面截面线性束，其每个成员 $Y_\xi$ 都有性质 $\mathcal P$。\footnote{法稿此处引用的节号没有定稿；英稿的“29”是其脚注号，不是所指节号。}

从编写角度看，本节的结果细化了前面若干断言。究竟应单列一节（或一条）补充这些细化，还是在前文逐次给出更精确的版本？\emph{Redactor decidetur}（作者的拉丁语附记，意在请编者决定）。
\section*{8.2\quad 维数跳跃和模的正则性}
设 $f:X\to P$ 是前文的态射。记 $Z$ 为使 $\dim Y_\xi>\dim X-1$ 的参数 $\xi\in P^\vee$ 所成的集合。假设对 $X$ 的每个不可约分支 $X_i$，均有 $\dim f(X_i)>0$，则 $Z$ 的闭包在 $P^\vee$ 中余维至少为 $2$。\footnote{法31的右缘裁掉了严格不等号后的阈值，英21也存疑；同段明写相关分支的正维条件，且这一条件对应下文的线性方程证明。此处按正维复原。法31首段把余维写作“2”，英21亦然；一般只能断言至少2。法31同段的模断言则明写“$\ge2$”。像张成一个平面时，包含该像的超平面参数余维已经是3；不存在此类超平面时例外集为空。}
作者指出，可由2.1、2.2（它们保证通用关联概形 $Y$ 的每个不可约分支支配 $P^\vee$）及 $Y\to P^\vee$ 的维数理论得出这一点。

作为推论，可从 $X$ 的闭子集 $T$ 出发，考察其在各个 $Y_\xi$ 中的逆像 $T_\xi$ 的维数。还可以把 $Z$ 定义为如下参数的集合：$T_{k(\xi)}$ 有一个不可约分支，它在 $Y_\xi$ 上的截面维数过大。仍假设 $T$ 的每个不可约分支的像均为正维。\footnote{法31的措辞是截面维数“过大”；英21误导为“取最大维数”。此处按相对于正常截面维数的跳跃理解。}

最后，作者把该方向上更精细的结论写为模的版本，并指向第一个断言的不可约情形和2.7。设 $\mathcal F$ 是 $X$ 上的凝聚模。若对 $\mathcal F$ 的每一个关联素循环 $T$ 都有 $\dim f(T)>0$，则使 $\varphi_\xi$ 不是 $\mathcal F_{k(\xi)}$-正则的参数 $\xi$ 组成可构造集，闭包的余维至少为 $2$。$\varphi_\xi$ 的记号沿用第5节。作者提议把它当作主命题，把前述断言列为推论；实际证明可先经由其中一个推论。

用前述记号，若 $\xi\in P^\vee\setminus Z$（这里 $Z$ 是上文模的非正则参数集），则对每个 $y\in Y_\xi$，取截面之前的模和取截面后的模 $\mathcal G_\xi$ 在 $y$ 的余深度相同。因此，若 $\coprof\mathcal F\le n$，则 $\coprof\mathcal G_\xi\le n$；特别地，若 $\mathcal F$ 是 Cohen--Macaulay 模，则 $\mathcal G_\xi$ 也是。最后，若 $\mathcal F$ 满足 $(S_k)$，则在上述参数范围内 $\mathcal G_\xi$ 满足 $(S_{k-1})$，参见第四章。\footnote{附记在法32开头。法稿等式左端的字形近似 $B_{k(\xi)}$，英22作 $G_{k(\xi)}$，下一行又将截面模一度误排为 $T_\xi$。按本段及第5节的正则截面语义，正文明确写出取截面前后的两个模；不将字形差异当作新对象。}
\section*{校记：两个划销片段}
法30下部旧8.2从 $X$ 的一个局部闭子集 $T$ 出发，考察 $T_\xi$ 比泛纤维维数大之处：当 $\dim\overline{f(T)}>0$ 时，所写的比较目标是截后仍为 $\dim T$；当 $\dim\overline{f(T)}=0$ 时，则改为考察截面非空。旧稿随后直接写例外集余维2，并提议推广到模的支集和 $T$ 的几何不可约分支。末句在指到某个不可约分支时截断。此段整块划销；正稿由法31页首另起，不能将它的结论倒接到旧稿。尤其若 $f$ 将非空 $T$ 映到 $k$-有理点，例外恰为过该点的参数超平面，余维仅为1。

法31页末另有整块框划的“T、$T'$ 的另一个推论”。可辨读出它比较两个闭集、提到2.4、再试图约化为 $T$ 不可约并考察余维2的闭例外集；右缘和底部裁切使关系符号及结句不足以可靠复原。这段不纳入正稿，不把可辨关键词拼成新命题。
\section*{独立核算：为什么是至少二}
\textbf{一、参数直线。}
设 $r=\dim P^\vee\ge2$，令 $W=\overline{Z_{\mathcal P}}$，且 $\dim W\le r-2$。直线参数空间 $G=\Gr(2,r+1)$ 的维数为 $2r-2$。在关联族 $\{(w,D):w\in W,\ w\in D\}$ 中，给定 $w$ 的直线族为 $\mathbb P^{r-1}$，故关联族维数至多 $2r-3$。其在 $G$ 中的像是闭的（投影固有），且不是全体。补集是定义在 $k$ 上的非空开集，必和稠密的仿射大胞相交。无限域上的非空仿射开集有 $k$-有理点，故有一条定义在 $k$ 上的直线在几何点意义下也避开 $W$。若 $r=1$，余维条件迫使 $W$ 为空，唯一的直线是全空间；$r=0$ 不存在直线。

\textbf{二、包含一个正维像的超平面。}
过一个点只给出一道线性条件；包含一个张成 $s$ 维射影线性空间的像则给出 $s+1$ 道独立条件。因为这些超平面的线性式需在一个 $s+1$ 维向量子空间上全为零，其非零解的射影化是余维 $s+1$ 的线性子空间，或为空。正维像有 $s\ge1$。在代数闭包上，有限型概形的不可约分支有限；原分支换底后的几何分支均支配原分支，且其像维数仍为正（可先在像的泛点考察换底再用维数不变）。所以有限个这类例外参数子空间的并，余维仍至少为2。该并在一个有限域扩张上定义且作为几何集合对基域共轭稳定；取各参数子空间的有限映射闭像可得到同一例外闭集在 $k$ 上的描述。

若截面不包含一个不可约分支，其交为空，或由非零局部方程截成余维1。于是 $\dim Y_\xi>\dim X-1$ 只会在有最大维分支被包含时发生；对各分支分别要求降维也同样有上面的余维下界。此论证不会把所有维数跳跃误判成恰余维2。

\newpage
\textbf{三、正则性与余深度。}
局部乘法非单射当且仅当方程落入有限模的某个关联素理想。扩域后的关联点投影到原来的关联点；在其闭包上仍可用上一段的正维像判据。相关的换底性质可见
\href{https://stacks.math.columbia.edu/tag/05EW}{Stacks Project，31.6}。
将 $\mathcal F$ 从 $k$ 扩到 $k(\xi)$ 时，支集维数和深度的局部增加量相同：局部换底是平坦的，纤维是有限型剩余域的扩域，局部为 Cohen--Macaulay；其维数和深度相等。由平坦局部换底的维数与深度公式，余深度不增加，$(S_k)$ 也保持。参见
\href{https://stacks.math.columbia.edu/tag/045O}{Stacks Project，33.13 的扩域论证}。
在截面上取一点，将截前模记作 $M$，正则方程记作 $h$，截后模为 $M/hM$。支集局部维数和深度各减少 $1$，因此
\[
 \dim\Supp(M/hM)-\depth(M/hM)=\dim\Supp M-\depth M.
\]
这里使用有限模的正则元深度公式，参见
\href{https://stacks.math.columbia.edu/tag/00LE}{Stacks Project，10.72}。
因为 $(S_k)$ 的局部条件是 $\depth M\ge\min(k,\dim\Supp M)$，取商后得到 $\depth(M/hM)\ge\min(k-1,\dim\Supp(M/hM))$。这一计算并不证明 $(S_k)$ 原样保留；下一条8.3 正要考察其失败可否发生在余维1的参数上。

\section*{底本与下一步}
底本：法文扫描、\href{https://webusers.imj-prg.fr/~leila.schneps/grothendieckcircle/EGA/EGA5trans.pdf}{英编辑稿}，以及包内原页对照。续自法32的8.3，先核其起笔划销与手写改正，再核法33是否接8.4。正式印刷本第四章第一分册物理221/223为另一进度线。
\end{document}

